\subsection*{Note (1)}
The first later note, on p.~52, explicitly retracts the judgment ``hors
de portée'' used in \S2, p.~5, and again in \S6, p.~42. Grothendieck now
calls the expression hasty and unfounded, treating it as a report on the
speaker's ideas and engagement rather than a measure of a problem's
mathematical possibilities. The joke about La Palisse accompanies his
references to the conjectures of Mordell, Tate, and Chafarévitch. He then
describes a change occurring in the days after he wrote the main report:
renewed contact with the earlier questions has brought a further property
of the outer Galois action on the fundamental group of an algebraic curve
to his attention. He presents this as progress toward the algebraic
characterization sought in \S2 and places it, together with the
``conjecture fondamentale'' of \S3, at the center of his anabelian
interests. The note gives neither that new property's statement nor a
proof; its function here is to record the reassessment and the renewed
direction of work.
(\href{https://grothendieck.umontpellier.fr/web/viewer.html?file=../119.pdf\#page=55}{Copy 1, p.~52, site 54};
\href{https://grothendieck.umontpellier.fr/web/viewer.html?file=../119.pdf\#page=8}{first call on p.~5, site 7};
\href{https://grothendieck.umontpellier.fr/web/viewer.html?file=../119.pdf\#page=45}{second call on p.~42, site 44}.)
